\chapter{Endliche Netze, Kompaktheit und Separabilität}
\label{b47guide:ch:compactness}

\section{Totale Beschränktheit}

\begin{definition}[Endliches Netz und totale Beschränktheit]
\label{b47guide:def:totally-bounded}
Für $A\subseteq X$ und $\varepsilon>0$ heißt eine endliche Menge
$F\subseteq A$ ein \emph{$\varepsilon$-Netz für $A$}, wenn
\[
 A\subseteq\bigcup_{a\in F}B_\varepsilon(a).
\]
$A$ heißt \emph{total beschränkt}, wenn für jedes $\varepsilon>0$ ein
solches Netz existiert. Für $A=\varnothing$ ist das leere Netz zugelassen.
\end{definition}

\begin{theorem}[Vererbung totaler Beschränktheit]
Jede endliche Menge ist total beschränkt. Total beschränkte Mengen sind
beschränkt. Teilmengen und Abschlüsse total beschränkter Mengen sind
wieder total beschränkt; ebenso endliche Vereinigungen solcher Mengen.
Insbesondere ist $A$ genau dann total beschränkt, wenn $\overline A$
total beschränkt ist.
\end{theorem}

\begin{theorem}[Cauchy-Teilfolgenkriterium]
\label{b47guide:thm:tb-cauchy}
$A\subseteq X$ ist genau dann total beschränkt, wenn jede Folge in $A$
eine Cauchy-Teilfolge besitzt. \cite{project,baire-notes}
\end{theorem}

\begin{theorem}[Getrennte Punkte als Gegenkriterium]
Eine Menge $A\subseteq X$ ist genau dann nicht total beschränkt, wenn es
ein $\varepsilon>0$ und eine Folge $(a_n)$ in $A$ gibt mit
\[
 d(a_n,a_m)\geq\varepsilon\qquad(n\ne m).
\]
\end{theorem}

\section{Charakterisierungen der Kompaktheit}

\begin{definition}[Folgenkompaktheit und Kompaktheit]
\label{b47guide:def:compact}
$K\subseteq X$ heißt \emph{folgenkompakt}, wenn jede Folge in $K$ eine
Teilfolge besitzt, die gegen einen Punkt von $K$ konvergiert.
Eine \emph{offene Überdeckung} von $K$ ist eine Familie von im Teilraum
$K$ offenen Mengen mit Vereinigung $K$. $K$ heißt \emph{kompakt}, wenn
jede solche Überdeckung eine endliche Teilüberdeckung besitzt.
Die leere Teilüberdeckung ist zugelassen.
\end{definition}

\begin{theorem}[Hauptcharakterisierung metrischer Kompaktheit]
\label{b47guide:thm:compact-equivalence}
Für jede Teilmenge $K$ eines metrischen Raums sind äquivalent:
\begin{enumerate}
\item $K$ ist kompakt.
\item $K$ ist folgenkompakt.
\item $(K,d_K)$ ist vollständig und $K$ ist total beschränkt.
\item Jede unendliche Teilmenge von $K$ besitzt einen Häufungspunkt in $K$.
\end{enumerate}
Die Vollständigkeit in Punkt~3 betrifft $K$ selbst, nicht den umgebenden
Raum. \cite{project,metric-notes,baire-notes}
\end{theorem}

\begin{theorem}[Kompakte Teilmengen und ihre abgeschlossenen Teile]
Kompakte Teilmengen metrischer Räume sind abgeschlossen und beschränkt.
Ist $K$ kompakt und $A\subseteq K$, so ist $A$ genau dann kompakt, wenn
$A$ im Teilraum $K$ abgeschlossen ist. Endliche Vereinigungen kompakter
Teilmengen sind kompakt.
\end{theorem}

\begin{definition}[Relative Kompaktheit]
$A\subseteq X$ heißt \emph{relativ kompakt in $X$}, wenn der Abschluss
$\overline A^{\,X}$ kompakt ist. Der umgebende Raum gehört zu dieser
Begriffsbestimmung.
\end{definition}

\begin{theorem}[Relative Kompaktheit in vollständigen Räumen]
\label{b47guide:thm:relative-compact}
Ist $X$ vollständig, so gilt für $A\subseteq X$:
\[
 A\text{ relativ kompakt in }X
 \quad\Longleftrightarrow\quad
 A\text{ total beschränkt}.
\]
In einem beliebigen metrischen Umraum gilt die Äquivalenz, wenn rechts
zusätzlich die Vollständigkeit von $\overline A^{\,X}$ gefordert wird.
\end{theorem}

\begin{theorem}[Lebesgue-Zahl einer offenen Überdeckung]
\label{b47guide:thm:lebesgue}
Ist $K$ kompakt und $\mathcal U$ eine offene Überdeckung von $K$, so
existiert ein $\lambda>0$ mit folgender Eigenschaft: Für jedes $x\in K$
ist die Teilraumkugel $B_\lambda^K(x)$ in einem Element von $\mathcal U$
enthalten. Insbesondere liegt jede nichtleere Teilmenge von $K$ mit
Durchmesser kleiner als $\lambda$ in einem Überdeckungselement.
\cite{metric-notes}
\end{theorem}

\begin{theorem}[Nächste Punkte und positive Abstände]
Ist $\varnothing\ne K\subseteq X$ kompakt, so besitzt jeder Punkt $x\in X$
einen nächsten Punkt in $K$. Ist außerdem $F\subseteq X$ nichtleer,
abgeschlossen und zu $K$ disjunkt, so gilt
\[
 d(K,F):=\inf\{d(k,f):k\in K,\ f\in F\}>0.
\]
Ist auch $F$ kompakt, wird dieses Infimum von einem Punktepaar angenommen.
Nächste Punkte müssen nicht eindeutig sein.
\end{theorem}

\section{Separabilität und abzählbare Kugelbeschreibung}

\begin{definition}[Separabilität und Basis]
Ein metrischer Raum heißt \emph{separabel}, wenn er eine höchstens
abzählbare dichte Teilmenge besitzt. Eine Familie $\mathcal B$ offener
Mengen heißt \emph{Basis}, wenn zu jeder offenen Menge $U$ und jedem
$x\in U$ ein $B\in\mathcal B$ mit $x\in B\subseteq U$ existiert.
\end{definition}

\begin{theorem}[Abzählbare Beschreibung metrischer Räume]
\label{b47guide:thm:separable}
Für einen metrischen Raum $X$ sind äquivalent:
\begin{enumerate}
\item $X$ ist separabel.
\item $X$ besitzt eine höchstens abzählbare Basis.
\item Jede offene Überdeckung von $X$ besitzt eine höchstens abzählbare
Teilüberdeckung.
\end{enumerate}
Ist $D\subseteq X$ höchstens abzählbar und dicht, so bilden die Kugeln
mit Mittelpunkten aus $D$ und positiven rationalen Radien eine solche Basis.
\cite{countable-notes}
\end{theorem}

\begin{theorem}[Vererbung der Separabilität]
Jeder Teilraum eines separablen metrischen Raums ist separabel. Jeder
total beschränkte metrische Raum ist separabel, insbesondere jeder
kompakte metrische Raum. Für $A\subseteq X$ sind $A$ und $\overline A$
entweder beide separabel oder beide nicht separabel.
\end{theorem}